\begin{abstract}
Autonomous navigation in dynamic platformer environments presents a classic reinforcement learning challenge characterized by delayed credit assignment and perceptual aliasing. This paper evaluates a tabular Q-learning agent deployed in a custom 2D graveyard platformer entitled \textit{Ho{\'a}'s Midnight Escape}, featuring a compact 144-state discretization space and a multi-stage curriculum progression. Training is governed by an asymmetric progress-shaped reward function balancing forward momentum and pit avoidance. While team members contributed environment physics, state discretization, and training orchestration, this paper focuses rigorously on empirical evaluation and statistical verification across 1,200 training episodes across three complexity tiers. The agent achieves a 99.0\% win rate on Level 1, 93.0\% on Level 2, and 88.0\% on Level 3 in the converged regime. Two-sample Welch's $t$-tests confirm statistically significant reward improvements from early exploration to convergence ($t = 23.41, 22.77, 22.90$ for Levels 1--3, all $p < 0.001$), accompanied by substantial reductions in navigation steps and collision rates. We present comprehensive learning curves, reward profiling benchmarks, playstyle clustering matrices, and failure mode decompositions demonstrating the efficacy of curriculum-based policy transfer in discrete state representations.

\keywords{Reinforcement Learning \and Tabular Q-Learning \and Curriculum Learning \and State Discretization \and Empirical Evaluation \and Statistical Hypothesis Testing.}
\end{abstract}
